\documentclass[conference]{IEEEtran}
\IEEEoverridecommandlockouts
\usepackage{cite}
\usepackage{amsmath,amssymb,amsfonts}
\usepackage{algorithmic}
\usepackage{graphicx}
\usepackage{textcomp}
\usepackage{xcolor}
\usepackage[brazil]{babel}

\def\BibTeX{{\rm B\kern-.05em{\sc i\kern-.025em b}\kern-.08em
    T\kern-.1667em\lower.7ex\hbox{E}\kern-.125emX}}  

\begin{document}

\title{Resenha sobre o Padrão de Projeto Facade}

\author{
\IEEEauthorblockN{Artur Bomtempo Colen}
\IEEEauthorblockA{
Engenharia de Software \\
Pontifícia Universidade Católica de Minas Gerais (PUC Minas) \\
Belo Horizonte, Brasil \\
abcolen@sga.pucminas.br}
}

\maketitle

\begin{abstract}

Este trabalho apresenta uma análise do material sobre o padrão de projeto \textit{Facade}, disponível no site Refactoring Guru. O conteúdo descreve um padrão estrutural cujo objetivo é fornecer uma interface simplificada para um subsistema complexo, reduzindo o acoplamento entre componentes e facilitando a utilização de funcionalidades pelo cliente. O material aborda o problema da complexidade em sistemas compostos por múltiplas classes, apresenta a solução por meio de uma fachada e discute aspectos como estrutura, aplicabilidade, vantagens e limitações. Esta resenha tem como objetivo explicar os principais conceitos apresentados, demonstrar o entendimento sobre o padrão e discutir sua aplicação prática no desenvolvimento de software.

\end{abstract}

\begin{IEEEkeywords}
engenharia de software, design patterns, facade, arquitetura de software, acoplamento, complexidade
\end{IEEEkeywords}

\section{Introdução}

No desenvolvimento de software, especialmente em sistemas de médio e grande porte, é comum lidar com múltiplas classes interdependentes. Essa complexidade tende a crescer com o tempo, principalmente quando o sistema depende de bibliotecas externas ou frameworks que possuem diversas funcionalidades internas.

O material analisado apresenta o padrão de projeto \textit{Facade}, que tem como principal objetivo simplificar essa complexidade ao fornecer uma interface única para um conjunto de classes. Em vez de o cliente interagir diretamente com várias partes do sistema, ele passa a utilizar uma única classe que encapsula toda a lógica necessária.

O conteúdo é estruturado de forma didática, apresentando o problema, a solução, exemplos, estrutura e aplicação. Isso permite compreender não apenas o funcionamento do padrão, mas também o contexto em que ele deve ser utilizado.

\section{Problema da Complexidade}

O material destaca que, ao trabalhar com sistemas complexos, o desenvolvedor frequentemente precisa lidar com diversos objetos, dependências e regras de execução. Conforme ilustrado no conteúdo, é necessário inicializar classes, controlar interações e garantir a ordem correta de chamadas de métodos.

Esse cenário gera um alto nível de acoplamento entre o código cliente e o subsistema. Como consequência, qualquer alteração na implementação interna pode exigir modificações em várias partes do sistema.

Além disso, a complexidade prejudica a legibilidade e a manutenção do código, pois o fluxo de execução fica distribuído entre várias classes. Isso torna o sistema mais difícil de entender, especialmente para novos desenvolvedores.

Outro ponto importante é que essa complexidade pode aumentar ainda mais quando se utiliza código de terceiros, já que o desenvolvedor não tem controle sobre sua estrutura interna.

\section{Solução: o padrão Facade}

Como solução, o material apresenta o padrão \textit{Facade}, que consiste na criação de uma classe responsável por fornecer uma interface simplificada para o subsistema.

De acordo com o conteúdo, a fachada encapsula toda a complexidade interna e expõe apenas as funcionalidades essenciais para o cliente. Dessa forma, o código cliente se torna mais simples e independente da implementação interna.

A analogia apresentada no material compara a fachada a um operador de atendimento, que atua como intermediário entre o cliente e diferentes setores de uma empresa. O cliente não precisa conhecer os detalhes de cada setor, apenas interagir com o operador.

A estrutura do padrão mostra que o cliente se comunica exclusivamente com a fachada, enquanto o subsistema permanece responsável por executar as operações reais. As classes internas não precisam conhecer a existência da fachada.

O exemplo de pseudocódigo apresentado reforça esse conceito ao mostrar como uma única classe pode encapsular várias operações de um sistema de conversão de vídeo, simplificando a interação com um framework complexo.

\section{Vantagens e Limitações}

O padrão Facade apresenta diversas vantagens. A principal delas é a redução do acoplamento, já que o cliente deixa de depender diretamente de múltiplas classes do subsistema.

Outro benefício importante é a melhoria na organização do código, pois a lógica de interação com o sistema fica centralizada em um único ponto. Isso facilita tanto a manutenção quanto a evolução do sistema.

Além disso, o padrão permite substituir ou modificar o subsistema sem impactar o cliente, desde que a interface da fachada seja mantida.

No entanto, o material também destaca algumas limitações. Uma fachada pode se tornar excessivamente complexa se acumular muitas responsabilidades, tornando-se difícil de manter. Nesse caso, pode ser necessário dividir a fachada em múltiplas classes.

Outro ponto é que o padrão não reduz a complexidade interna do sistema, apenas a esconde, o que pode ser uma desvantagem em alguns cenários.

\section{Aplicação Prática}

Na prática, o padrão Facade é amplamente utilizado no desenvolvimento de software moderno.

Um exemplo comum é a integração com APIs ou bibliotecas externas. Em vez de espalhar chamadas complexas pelo sistema, pode-se criar uma fachada que centraliza essas interações, facilitando o uso e a manutenção.

Outro exemplo é a organização de sistemas em camadas. Cada camada pode expor uma fachada que simplifica o acesso às suas funcionalidades, reduzindo o acoplamento entre componentes.

O padrão também é útil em sistemas que passam por mudanças frequentes. Ao centralizar a lógica de interação, torna-se mais fácil adaptar o sistema sem impactar o restante do código.

Durante projetos acadêmicos, a aplicação do Facade pode ajudar a organizar melhor o código e evitar dependências desnecessárias, tornando o sistema mais claro e estruturado.

\section{Conclusão}

O material do Refactoring Guru apresenta de forma clara e didática o padrão de projeto Facade, destacando sua importância na simplificação de sistemas complexos.

O principal entendimento obtido é que a criação de uma interface unificada permite reduzir o acoplamento e facilitar a interação com subsistemas, tornando o código mais organizado e compreensível.

Além disso, o padrão contribui para melhorar a manutenção e evolução do software, sendo uma solução prática para lidar com sistemas complexos.

Dessa forma, compreender e aplicar o padrão Facade é essencial para o desenvolvimento de sistemas mais bem estruturados, tanto em contextos acadêmicos quanto profissionais.

\begin{thebibliography}{00}

\bibitem{b1}
Refactoring Guru,
``Facade,''
Disponível em: https://refactoring.guru/pt-br/design-patterns/facade

\end{thebibliography}

\end{document}